\documentclass[a4paper,11pt]{article}

\usepackage[a4paper,margin=2cm]{geometry}
\usepackage[T1]{fontenc}
\usepackage[utf8]{inputenc}
\usepackage[ngerman,english]{babel}
\usepackage{lmodern}
\usepackage{microtype}
\usepackage{xcolor}
\usepackage{parskip}
\usepackage{enumitem}
\usepackage[most]{tcolorbox}
\usepackage{titlesec}
\usepackage[hidelinks]{hyperref}

\definecolor{HeaderBlueDark}{RGB}{70,110,170}
\definecolor{RowBlueLight}{RGB}{235,243,255}
\definecolor{RowBlueDark}{RGB}{220,233,250}
\definecolor{SoftGray}{RGB}{90,90,90}

\setlength{\parindent}{0pt}
\setlist[itemize]{leftmargin=1.4em,itemsep=0.35em,topsep=0.35em}
\linespread{1.06}

\titleformat{\section}[block]{\Large\bfseries\color{HeaderBlueDark}}{}{0pt}{}
\titlespacing{\section}{0pt}{1.3em}{0.8em}
\titleformat{\subsection}[block]{\large\bfseries\color{HeaderBlueDark}}{}{0pt}{}
\titlespacing{\subsection}{0pt}{1.0em}{0.6em}

\newtcolorbox{exbox}{enhanced,breakable,colback=RowBlueLight,colframe=RowBlueDark,boxrule=0.4pt,arc=1.5mm,left=1.2mm,right=1.2mm,top=1mm,bottom=1mm,title={Beispiele \textnormal{(Examples)}},fonttitle=\bfseries,colbacktitle=RowBlueDark,coltitle=black}
\NewDocumentEnvironment{grammarentry}{mm+b}{%
    \begin{tcolorbox}[enhanced,breakable,colback=white,colframe=HeaderBlueDark,boxrule=0.6pt,arc=1.5mm,left=1.5mm,right=1.5mm,top=1mm,bottom=1mm,colbacktitle=HeaderBlueDark,coltitle=white,fonttitle=\bfseries,title={#1\hfill\normalfont\footnotesize #2}]
    #3
    \end{tcolorbox}
}{}
\newcommand{\GermanText}[1]{\textbf{Deutsch: }#1\par}
\newcommand{\EnglishText}[1]{{\small\color{SoftGray}\textbf{English: }#1\par}}
\newcommand{\ExampleItem}[2]{\item \textbf{#1}\\{\small\color{SoftGray}#2}}

\title{Die Präposition \glqq bei\grqq\\\large The Preposition ``bei''}
\date{}

\begin{document}
\maketitle
\tableofcontents
\newpage

\section{Grundbedeutung -- Basic meaning}
\begin{grammarentry}{bei}{Präposition + Dativ / preposition + dative}
\GermanText{\textbf{bei} ist eine Präposition, die immer den \textbf{Dativ} verlangt. Sie kann einen Ort, eine Person, eine Institution, eine Tätigkeit, eine Situation, einen Zeitpunkt oder eine Bedingung bezeichnen.}
\EnglishText{\textbf{bei} is a preposition that always requires the \textbf{dative}. It can refer to a place, person, institution, activity, situation, time, or condition.}
\GermanText{Merksatz: \textbf{bei + Dativ}.}
\EnglishText{Rule to remember: \textbf{bei + dative}.}
\end{grammarentry}

\section{Kasus -- Case}
\begin{grammarentry}{Dativ}{immer Dativ / always dative}
\GermanText{Nach \textbf{bei} steht immer der Dativ: \textbf{bei dem Mann}, \textbf{bei der Frau}, \textbf{bei dem Kind}, \textbf{bei den Freunden}.}
\EnglishText{After \textbf{bei}, the dative is always used: \textbf{bei dem Mann}, \textbf{bei der Frau}, \textbf{bei dem Kind}, \textbf{bei den Freunden}.}
\GermanText{Im Plural erhält das Nomen im Dativ normalerweise zusätzlich ein \textbf{-n}, wenn die Pluralform nicht schon auf -n oder -s endet: \textbf{bei den Kindern}, \textbf{bei den Freunden}.}
\EnglishText{In the plural dative, the noun normally also gets an \textbf{-n} if the plural form does not already end in -n or -s: \textbf{bei den Kindern}, \textbf{bei den Freunden}.}
\end{grammarentry}

\section{Verschmelzung mit dem Artikel -- Contraction with the article}
\begin{grammarentry}{beim}{bei + dem}
\GermanText{\textbf{bei dem} wird sehr häufig zu \textbf{beim} zusammengezogen. Das gilt für maskuline und neutrale Nomen im Dativ Singular.}
\EnglishText{\textbf{bei dem} is very often contracted to \textbf{beim}. This applies to masculine and neuter nouns in the dative singular.}
\GermanText{\textbf{bei der}, \textbf{bei den} und \textbf{bei einer} werden nicht entsprechend zusammengezogen.}
\EnglishText{\textbf{bei der}, \textbf{bei den}, and \textbf{bei einer} do not have equivalent standard contractions.}
\end{grammarentry}
\begin{exbox}\begin{itemize}
\ExampleItem{Ich bin beim Arzt.}{I am at the doctor's.}
\ExampleItem{Wir treffen uns beim Bahnhof.}{We are meeting near the station.}
\ExampleItem{Sie ist bei der Ärztin.}{She is at the female doctor's.}
\end{itemize}\end{exbox}

\section{Bedeutung 1: bei einer Person -- Meaning 1: at a person's place}
\begin{grammarentry}{Person}{bei jemandem / at someone's place}
\GermanText{\textbf{bei} bezeichnet sehr häufig, dass man sich bei einer Person, in ihrem Haus, in ihrer Wohnung oder in ihrer unmittelbaren Umgebung befindet.}
\EnglishText{\textbf{bei} very often indicates that someone is at another person's place, in their house or apartment, or in their immediate surroundings.}
\end{grammarentry}
\begin{exbox}\begin{itemize}
\ExampleItem{Ich bin heute bei meiner Schwester.}{I am at my sister's today.}
\ExampleItem{Wir essen morgen bei Freunden.}{We are eating at friends' place tomorrow.}
\ExampleItem{Er wohnt noch bei seinen Eltern.}{He still lives with his parents.}
\end{itemize}\end{exbox}

\section{Bedeutung 2: Firma, Institution oder Organisation -- Meaning 2: company, institution, or organization}
\begin{grammarentry}{Institution}{Arbeitsplatz oder Zugehörigkeit / workplace or affiliation}
\GermanText{Mit \textbf{bei} sagt man oft, für welche Firma oder Institution jemand arbeitet oder zu welcher Organisation eine Person gehört.}
\EnglishText{\textbf{bei} is often used to say which company or institution someone works for or which organization someone belongs to.}
\end{grammarentry}
\begin{exbox}\begin{itemize}
\ExampleItem{Sie arbeitet bei Siemens.}{She works at Siemens.}
\ExampleItem{Er ist bei der Polizei.}{He is with the police.}
\ExampleItem{Ich habe einen Termin bei der Bank.}{I have an appointment at the bank.}
\end{itemize}\end{exbox}

\section{Bedeutung 3: räumliche Nähe -- Meaning 3: spatial proximity}
\begin{grammarentry}{Nähe}{nahe an / near}
\GermanText{\textbf{bei} kann bedeuten, dass sich etwas in der Nähe eines Ortes befindet. Es bedeutet dann ungefähr \textbf{in der Nähe von}.}
\EnglishText{\textbf{bei} can mean that something is located near a place. In this use it means approximately \textbf{near} or \textbf{in the vicinity of}.}
\GermanText{\textbf{bei} ist weniger genau als \textbf{an}. Es bezeichnet oft nur die allgemeine Nähe und nicht unbedingt direkten Kontakt oder eine genaue Position.}
\EnglishText{\textbf{bei} is less precise than \textbf{an}. It often indicates only general proximity, not necessarily direct contact or an exact position.}
\end{grammarentry}
\begin{exbox}\begin{itemize}
\ExampleItem{Potsdam liegt bei Berlin.}{Potsdam is near Berlin.}
\ExampleItem{Wir wohnen bei der Universität.}{We live near the university.}
\end{itemize}\end{exbox}

\section{Bedeutung 4: Tätigkeit oder Vorgang -- Meaning 4: activity or process}
\begin{grammarentry}{bei + substantivierte Tätigkeit}{während einer Tätigkeit / during an activity}
\GermanText{\textbf{bei} kann ausdrücken, dass etwas während einer Tätigkeit oder eines Vorgangs passiert. Sehr häufig steht danach ein substantivierter Infinitiv.}
\EnglishText{\textbf{bei} can express that something happens during an activity or process. A nominalized infinitive often follows it.}
\GermanText{Struktur: \textbf{beim + substantivierter Infinitiv}.}
\EnglishText{Structure: \textbf{beim + nominalized infinitive}.}
\end{grammarentry}
\begin{exbox}\begin{itemize}
\ExampleItem{Beim Essen spreche ich nicht gern über Arbeit.}{I do not like talking about work while eating.}
\ExampleItem{Beim Autofahren darf man nicht aufs Handy schauen.}{You must not look at your phone while driving.}
\ExampleItem{Ich höre beim Lernen klassische Musik.}{I listen to classical music while studying.}
\end{itemize}\end{exbox}

\section{Bedeutung 5: Situation oder Umstand -- Meaning 5: situation or circumstance}
\begin{grammarentry}{Umstand}{unter einer bestimmten Bedingung / under a certain circumstance}
\GermanText{\textbf{bei} kann einen Umstand oder eine Situation angeben, unter der etwas gilt oder passiert.}
\EnglishText{\textbf{bei} can specify a circumstance or situation under which something applies or happens.}
\end{grammarentry}
\begin{exbox}\begin{itemize}
\ExampleItem{Bei Regen bleiben wir zu Hause.}{When it rains / In case of rain, we stay at home.}
\ExampleItem{Bei Problemen ruf mich an.}{If there are problems, call me.}
\ExampleItem{Bei hoher Geschwindigkeit wird das Fahren gefährlich.}{At high speed, driving becomes dangerous.}
\end{itemize}\end{exbox}

\section{Bedeutung 6: Zeitpunkt oder Gelegenheit -- Meaning 6: time or occasion}
\begin{grammarentry}{Gelegenheit}{bei einem Ereignis / on an occasion}
\GermanText{\textbf{bei} kann bedeuten, dass etwas im Zusammenhang mit einem Ereignis oder zu einer bestimmten Gelegenheit geschieht.}
\EnglishText{\textbf{bei} can mean that something happens in connection with an event or on a particular occasion.}
\end{grammarentry}
\begin{exbox}\begin{itemize}
\ExampleItem{Bei unserem letzten Treffen war er sehr müde.}{At our last meeting, he was very tired.}
\ExampleItem{Bei der Prüfung darfst du kein Wörterbuch benutzen.}{During the exam, you may not use a dictionary.}
\end{itemize}\end{exbox}

\section{Bedeutung 7: körperlicher oder persönlicher Zustand -- Meaning 7: physical or personal condition}
\begin{grammarentry}{Zustand}{unter bestimmten Bedingungen / in a certain condition}
\GermanText{\textbf{bei} wird auch mit Zuständen, Temperaturen, Mengen oder anderen messbaren Bedingungen verwendet.}
\EnglishText{\textbf{bei} is also used with states, temperatures, quantities, or other measurable conditions.}
\end{grammarentry}
\begin{exbox}\begin{itemize}
\ExampleItem{Wasser gefriert bei null Grad Celsius.}{Water freezes at zero degrees Celsius.}
\ExampleItem{Bei dieser Hitze kann ich nicht arbeiten.}{I cannot work in this heat.}
\end{itemize}\end{exbox}

\section{bei mit Pronomen -- bei with pronouns}
\begin{grammarentry}{Dativpronomen}{bei mir, bei dir, bei ihm ...}
\GermanText{Da \textbf{bei} den Dativ verlangt, verwendet man die Dativformen der Personalpronomen: \textbf{bei mir, bei dir, bei ihm, bei ihr, bei uns, bei euch, bei ihnen, bei Ihnen}.}
\EnglishText{Because \textbf{bei} requires the dative, the dative forms of personal pronouns are used: \textbf{bei mir, bei dir, bei ihm, bei ihr, bei uns, bei euch, bei ihnen, bei Ihnen}.}
\end{grammarentry}
\begin{exbox}\begin{itemize}
\ExampleItem{Du kannst heute bei mir bleiben.}{You can stay at my place today.}
\ExampleItem{Ist alles bei dir in Ordnung?}{Is everything okay with you?}
\end{itemize}\end{exbox}

\section{Wo? und Wohin? -- Where? and Where to?}
\begin{grammarentry}{keine Wechselpräposition}{bei bleibt Dativ / bei stays dative}
\GermanText{\textbf{bei} ist \textbf{keine Wechselpräposition}. Auch wenn eine Bewegung gemeint ist, verwendet man nicht den Akkusativ. Für eine Bewegung zu einer Person oder Institution benutzt man meistens \textbf{zu + Dativ}.}
\EnglishText{\textbf{bei} is \textbf{not a two-way preposition}. Even when movement is involved, the accusative is not used. For movement toward a person or institution, German usually uses \textbf{zu + dative}.}
\GermanText{Ort: \textbf{Ich bin beim Arzt.} Bewegung/Ziel: \textbf{Ich gehe zum Arzt.}}
\EnglishText{Location: \textbf{I am at the doctor's.} Movement/destination: \textbf{I am going to the doctor.}}
\end{grammarentry}

\section{Unterschied zwischen bei und zu -- Difference between bei and zu}
\begin{grammarentry}{bei vs. zu}{Ort gegen Richtung / location versus direction}
\GermanText{\textbf{bei} beschreibt normalerweise, wo jemand oder etwas ist. \textbf{zu} beschreibt meistens eine Bewegung oder Richtung zu einer Person, Institution oder einem Ziel.}
\EnglishText{\textbf{bei} normally describes where someone or something is. \textbf{zu} usually describes movement or direction toward a person, institution, or destination.}
\end{grammarentry}
\begin{exbox}\begin{itemize}
\ExampleItem{Ich bin bei meiner Mutter.}{I am at my mother's place.}
\ExampleItem{Ich fahre zu meiner Mutter.}{I am going to my mother's place.}
\ExampleItem{Er ist bei der Arbeit.}{He is at work.}
\ExampleItem{Er fährt zur Arbeit.}{He is going to work.}
\end{itemize}\end{exbox}

\section{Unterschied zwischen bei und mit -- Difference between bei and mit}
\begin{grammarentry}{bei vs. mit}{Ort/Zugehörigkeit gegen Begleitung/Mittel / location-affiliation versus accompaniment-means}
\GermanText{\textbf{bei} bezeichnet oft einen Ort, eine Nähe oder eine Zugehörigkeit. \textbf{mit} bezeichnet meist Begleitung, ein Mittel oder ein Werkzeug.}
\EnglishText{\textbf{bei} often indicates location, proximity, or affiliation. \textbf{mit} usually indicates accompaniment, a means, or a tool.}
\end{grammarentry}
\begin{exbox}\begin{itemize}
\ExampleItem{Ich bin bei Anna.}{I am at Anna's place.}
\ExampleItem{Ich bin mit Anna.}{I am with Anna.}
\ExampleItem{Ich arbeite bei einer Bank.}{I work at a bank.}
\ExampleItem{Ich bezahle mit einer Bankkarte.}{I pay with a bank card.}
\end{itemize}\end{exbox}

\section{Unterschied zwischen bei und an -- Difference between bei and an}
\begin{grammarentry}{bei vs. an}{allgemeine Nähe gegen genauere Position / general proximity versus more precise position}
\GermanText{\textbf{bei} bezeichnet oft allgemeine Nähe. \textbf{an} bezeichnet häufig eine genauere Position an einem Rand, einer Fläche oder einem bestimmten Punkt.}
\EnglishText{\textbf{bei} often indicates general proximity. \textbf{an} often indicates a more precise position at an edge, surface, or specific point.}
\end{grammarentry}
\begin{exbox}\begin{itemize}
\ExampleItem{Wir wohnen bei der Universität.}{We live near the university.}
\ExampleItem{Wir warten an der Bushaltestelle.}{We are waiting at the bus stop.}
\end{itemize}\end{exbox}

\section{Pronominaladverb dabei -- Pronominal adverb dabei}
\begin{grammarentry}{dabei}{da + bei}
\GermanText{Für Sachen oder abstrakte Inhalte verwendet man oft das Pronominaladverb \textbf{dabei}. Es kann je nach Kontext \glqq bei dieser Sache\grqq, \glqq währenddessen\grqq oder \glqq anwesend\grqq bedeuten.}
\EnglishText{For things or abstract content, German often uses the pronominal adverb \textbf{dabei}. Depending on context, it can mean ``in connection with it,'' ``while doing so,'' or ``present / there.''}
\GermanText{Für Personen sagt man dagegen: \textbf{bei ihm, bei ihr, bei ihnen} usw.}
\EnglishText{For people, however, German says: \textbf{bei ihm, bei ihr, bei ihnen}, etc.}
\end{grammarentry}
\begin{exbox}\begin{itemize}
\ExampleItem{Ich war bei der Besprechung dabei.}{I was present at the meeting.}
\ExampleItem{Er kocht und hört dabei Musik.}{He cooks and listens to music while doing so.}
\end{itemize}\end{exbox}

\section{Frageform wobei -- Question form wobei}
\begin{grammarentry}{wobei}{wo + bei}
\GermanText{Wenn nach einer Sache, Tätigkeit oder Situation mit \textbf{bei} gefragt wird, verwendet man oft \textbf{wobei}. Für Personen verwendet man \textbf{bei wem}.}
\EnglishText{When asking about a thing, activity, or situation with \textbf{bei}, German often uses \textbf{wobei}. For people, it uses \textbf{bei wem}.}
\end{grammarentry}
\begin{exbox}\begin{itemize}
\ExampleItem{Wobei brauchst du Hilfe?}{What do you need help with?}
\ExampleItem{Bei wem wohnst du?}{Who are you staying with?}
\end{itemize}\end{exbox}

\section{Häufige feste Verbindungen -- Common fixed expressions}
\begin{grammarentry}{Kollokationen}{typische Verbindungen / common combinations}
\GermanText{Häufig sind: \textbf{bei der Arbeit}, \textbf{bei Bedarf}, \textbf{bei Gelegenheit}, \textbf{bei Problemen}, \textbf{bei Regen}, \textbf{bei Nacht}, \textbf{bei Tageslicht}, \textbf{bei jemandem zu Gast sein}, \textbf{bei jemandem bleiben}, \textbf{bei einer Firma arbeiten}.}
\EnglishText{Common expressions include: \textbf{at work}, \textbf{if needed}, \textbf{when the opportunity arises}, \textbf{in case of problems}, \textbf{in the rain}, \textbf{at night}, \textbf{in daylight}, \textbf{to be a guest at someone's place}, \textbf{to stay with someone}, \textbf{to work at a company}.}
\end{grammarentry}

\section{Häufige Fehler -- Common mistakes}
\begin{grammarentry}{Fehler 1}{Akkusativ statt Dativ / accusative instead of dative}
\GermanText{Falsch: \textbf{bei den Arzt}. Richtig: \textbf{bei dem Arzt / beim Arzt}.}
\EnglishText{Wrong: \textbf{bei den Arzt}. Correct: \textbf{bei dem Arzt / beim Arzt}.}
\end{grammarentry}
\begin{grammarentry}{Fehler 2}{bei für eine Richtung / bei for direction}
\GermanText{Falsch für eine Bewegung: \textbf{Ich gehe bei meinem Freund}. Richtig: \textbf{Ich gehe zu meinem Freund}.}
\EnglishText{Wrong for movement: \textbf{Ich gehe bei meinem Freund}. Correct: \textbf{Ich gehe zu meinem Freund}.}
\end{grammarentry}
\begin{grammarentry}{Fehler 3}{bei und mit verwechseln / confusing bei and mit}
\GermanText{\textbf{Ich bin bei meinem Freund} bedeutet normalerweise: Ich bin an seinem Ort. \textbf{Ich bin mit meinem Freund} bedeutet: Er begleitet mich oder wir sind zusammen.}
\EnglishText{\textbf{Ich bin bei meinem Freund} normally means: I am at his place. \textbf{Ich bin mit meinem Freund} means: he is accompanying me or we are together.}
\end{grammarentry}

\section{Zusammenfassung -- Summary}
\begin{grammarentry}{Merksätze}{kurz und wichtig / short and important}
\GermanText{1. \textbf{bei + Dativ} -- immer.}
\EnglishText{1. \textbf{bei + dative} -- always.}
\GermanText{2. \textbf{beim = bei dem}.}
\EnglishText{2. \textbf{beim = bei dem}.}
\GermanText{3. \textbf{bei} kann Ort, Nähe, Person, Institution, Tätigkeit, Situation, Zeitpunkt oder Bedingung ausdrücken.}
\EnglishText{3. \textbf{bei} can express location, proximity, person, institution, activity, situation, time, or condition.}
\GermanText{4. Für Bewegung zu einer Person oder Institution benutzt man meistens \textbf{zu + Dativ}: \textbf{zum Arzt}, \textbf{zur Bank}, \textbf{zu meinen Eltern}.}
\EnglishText{4. For movement toward a person or institution, German usually uses \textbf{zu + dative}: \textbf{zum Arzt}, \textbf{zur Bank}, \textbf{zu meinen Eltern}.}
\GermanText{5. Für Dinge: \textbf{dabei / wobei}; für Personen: \textbf{bei ihm / bei wem}.}
\EnglishText{5. For things: \textbf{dabei / wobei}; for people: \textbf{bei ihm / bei wem}.}
\end{grammarentry}

\end{document}
